\scanpage{308}
\mspage{7}
quatre avatars possibles
\[
 K^*\quad K_*\quad K^{\cdot}\quad K_{\cdot}
\]
liés par les relations ci-dessus
\[
 \left\{\begin{array}{ll}
 K^*=DP^*(K^{\cdot})&K^{\cdot}=ND^{\cdot}(K^*)\\
 K_*=DP_*(K_{\cdot})&K_{\cdot}=ND_{\cdot}(K_*)
 \end{array}\right.
\]
plus les relations
\[
 \left\{\begin{array}{ll}
 K^*=(K_*)^\vee,&K_*=(K^*)^\vee\\
 K^{\cdot}=(K_{\cdot})^\vee,&K_{\cdot}=(K^{\cdot})^\vee.
 \end{array}\right.
\]
\deleted{\unsure{(et leur signification}}

Ces \deleted{$\unsure{K^*=}$} notations et remarques (i.e. équivalences de
catégories) s'étendent aux objets \deleted{\unsure{simpliciaux}}
\inserted{dépendant de} plusieurs indices, p.ex \deleted{$\unsure{K^*}$}
une même lettre $K$, désignant \inserted{\unsure{initialemt|initialement}}
un objet ${K^*}_*$ de $(A_*)^*$ \inserted{disons}, désigne également,
\deleted{\unsure{un objet}} quatre avatars distincts
\[
 \left\{\begin{array}{ll}
 {K^*}_*\in(A_*)^*,&({K^{\cdot}}_*)\in(A^{\deleted{\unsure{*}}}_*)^{\cdot}\\
 {K^{\cdot}}_{\cdot}\in(A_{\cdot})^{\cdot},&({K^*}_{\cdot})\in(A_{\cdot})^*
 \end{array}\right.
\]
avec des notations
\[
 \begin{array}{ll}
 {K^{[n]}}_{[m]},&{K^n}_{[m]}\\
 {K^n}_m,&{K^{[n]}}_m
 \end{array}
\]
pour les composantes. De même, on aura, si on part p.ex. de
$K^{**}\in(A^*)^*$
\[
 \left\{\begin{array}{ll}
 K^{**}\in(A^*)^*,&K^{\unsure{\cdot}*}\in(A^{\cdot})^*\\
 K^{\cdot\cdot}\in(A^{\cdot})^{\cdot},&K^{*\cdot}\in(A^*)^{\cdot}
 \end{array}\right.
\]
\[
 \left\{\begin{array}{ll}
 K^{[n][m]},&K^{n[m]}\\
 K^{nm},&K^{[n]m}
 \end{array}\right.
\]
On notera qu'un objet de $(A_*)^{\cdot}$ est un complexe de cochaînes
\inserted{$n\longmapsto({K_*}^{n},d^n)$}
dans la catég. $A_*$ des objets $\tfrac12$ simpliciaux de $A$,
\margin{309}
